\section{Computation}
\label{sec:computation}

We solve the partially penalized smoothed-rank objective in the original coefficient coordinates. For a fixed training sample and $(K,h,\lambda)$ with $h,\lambda>0$, write $\widehat{\bm{W}}=(\widehat{\bm{A}},\widehat{\bm{U}})$, $\widehat{\bm{A}}=(\bm{Z}_c,\widehat{\bm{F}})$, $q=d+K$, $P=q+p$, and $\bm{\vartheta}=(\bm{a}^\top,\bm{\beta}^\top)^\top$. The objective is
\begin{equation}
\begin{aligned}
Q(\bm\vartheta)&=f_h(\bm\vartheta)+g_\lambda(\bm\vartheta),\\
g_\lambda(\bm\vartheta)&=\lambda\|\bm\beta\|_1,\\
f_h(\bm\vartheta)&=\frac1N\sum_{i\ne j}\Delta_i\ell_h(r_{ij}),\\
r_{ij}&=Y_j-Y_i+(\widehat{\bm W}_i-\widehat{\bm W}_j)^\top\bm\vartheta,\\
N&=n(n-1).
\end{aligned}
\label{eq:comp-objective}
\end{equation}
where $Y_i$ is already on the log-time scale and $\ell_h(r)=r\Phi(r/h)+h\phi(r/h)$. All arrays are finite, $n\ge2$, and $\Delta_i\in\{0,1\}$. PCA and preprocessing remain fixed during a solve.

The reference algorithm is monotone proximal gradient with backtracking. The empirical implementation combines acceleration and qualified second-order refinement. These algorithms specialize composite-convex optimization \citep{BeckTeboulle2009,ParikhBoyd2014} to the pairwise loss and unpenalized block in~\eqref{eq:comp-objective}. Scalable alternatives for the unsmoothed AFT criterion include the prox-linear alternating direction method of multipliers (ADMM) algorithm of \citet{SuderMolstad2022}.

\subsection{Blocked pairwise evaluation}
\label{subsec:comp-evaluation}
Compute $\bm{e}=\bm Y-\widehat{\bm{W}}\bm{\vartheta}$. For $i\ne j$, put $V_{ij}=\Delta_i\Phi\{(e_j-e_i)/h\}$, and set $V_{ii}=0$. The gradient is
\begin{equation}
 c_i=\sum_jV_{ij}-\sum_jV_{ji},\qquad
 \nabla f_h(\bm{\vartheta})=\widehat{\bm{W}}^\top \bm{c}/N.
 \label{eq:comp-gradient}
\end{equation}
The coefficients satisfy $\mathbf 1_n^\top \bm{c}=0$. Event patients anchor rows $i$, all other patients contribute as $j$, and the diagonal is excluded from the ordered-pair sum.

For an event-row block $\mathcal I$, form the $|\mathcal I|\times n$ residual differences, accumulate the row and column contributions to $\bm{c}$, and sum the associated losses. This evaluates the complete ordered-pair criterion exactly. With $m=\sum_i\Delta_i$ events and block size $b$, one objective--gradient pass uses $O(nP+mn)$ arithmetic and $O(bn+n+P)$ temporary storage in addition to the stored design. Line-search evaluations contribute further work.

The residual matrix need not be stored either. For any $\bm{v}=(\bm{v}_\theta^\top,\bm{v}_\gamma^\top,\bm{v}_\beta^\top)^\top$ and $\bm{c}\in\mathbb R^n$, the identities
\begin{equation}
 \begin{split}
 \widehat{\bm{W}}\bm{v}&=\bm{Z}_c\bm{v}_\theta+\widehat{\bm{F}}(\bm{v}_\gamma-\widehat{\bm{B}}^\top \bm{v}_\beta)+\bm{X}_c\bm{v}_\beta,\\
 \widehat{\bm{W}}^\top \bm{c}&=
 \begin{pmatrix}\bm{Z}_c^\top \bm{c}\\ \widehat{\bm{F}}^\top \bm{c}\\
 \bm{X}_c^\top \bm{c}-\widehat{\bm{B}}(\widehat{\bm{F}}^\top \bm{c})\end{pmatrix}
 \end{split}
 \label{eq:comp-implicit-W}
\end{equation}
allow storage of $\bm{X}_c$ and the low-rank factors in place of a duplicate residual matrix. Euclidean updates remain in the original $(\bm{a},\bm{\beta})$ coordinates.

\subsection{Smoothness and a partially penalized proximal step}
\label{subsec:comp-smoothness}
Let $\bm{\mathscr L}_\Delta$ be the graph Laplacian with edge weights $\Delta_i+\Delta_j$ for $i<j$. Explicitly,
\begin{equation}
 \begin{gathered}
\bm{\mathscr L}_\Delta=n\operatorname{diag}(\bm\Delta)+m\bm{I}_n
             -\bm\Delta\mathbf 1_n^\top-\mathbf 1_n\bm\Delta^\top,
 \\
 \mathscr B_h=\frac{\|\widehat{\bm{W}}^\top\bm{\mathscr L}_\Delta\widehat{\bm{W}}\|_{\rm op}}
                    {Nh\sqrt{2\pi}}.
 \end{gathered}
 \label{eq:comp-lipschitz}
\end{equation}
Here $\bm\Delta=(\Delta_1,\ldots,\Delta_n)^\top$. Pair grouping and $\phi\le\phi(0)$ imply
\begin{equation}
 \begin{gathered}
0\preceq\nabla^2f_h(\bm{\vartheta})
 \preceq\frac{\widehat{\bm{W}}^\top\bm{\mathscr L}_\Delta\widehat{\bm{W}}}{Nh\sqrt{2\pi}},
 \\
 \|\nabla f_h(\bm{x})-\nabla f_h(\bm{y})\|_2\le\mathscr B_h\|\bm{x}-\bm{y}\|_2.
 \end{gathered}
 \label{eq:comp-hessian-bound}
\end{equation}
This global Lipschitz bound controls numerical step selection. Since $\widehat{\bm{W}}$ is centered,
\[\begin{gathered}
 \widehat{\bm{W}}^\top\bm{\mathscr L}_\Delta\widehat{\bm{W}}
 =n\widehat{\bm{W}}^\top\operatorname{diag}(\bm\Delta)\widehat{\bm{W}}+m\widehat{\bm{W}}^\top\widehat{\bm{W}},
 \\
 \mathscr B_h\le\frac{(n+m)\|\widehat{\bm{W}}\|_{\rm op}^2}{Nh\sqrt{2\pi}}.
\end{gathered}\]
The Frobenius norm gives a computable upper bound in place of the operator norm. Backtracking also permits step selection without evaluating these norms or forming a full Hessian. The factor $1/h$ governs numerical smoothness, separately from the stochastic scale in Section~\ref{sec:theory}.

For $t>0$, define the soft-thresholding map by $\mathsf S_\tau(z)_j=\operatorname{sign}(z_j)(|z_j|-\tau)_+$. The exact proximal-gradient candidate is
\begin{equation}
 \begin{split}
 T_t(\bm{x})&=\operatorname{prox}_{t g_\lambda}\{\bm{x}-t\nabla f_h(\bm{x})\},\\
 [T_t(\bm{x})]_a&=\bm{x}_a-t\nabla_a f_h(\bm{x}),\qquad
 [T_t(\bm{x})]_\beta=\mathsf S_{t\lambda}\{\bm{x}_\beta-t\nabla_\beta f_h(\bm{x})\}.
 \end{split}
 \label{eq:comp-prox}
\end{equation}
The proximal map leaves the nuisance block unpenalized. Rescaled residual coordinates use the corresponding weighted threshold from Section~\ref{sec:estimation}.

\subsection{Backtracking algorithm and deterministic convergence}
\label{subsec:comp-algorithm}
Choose $t_{\max}>0$ and $\zeta,\omega\in(0,1)$. From $\bm{x}=\bm{x}^{(r)}$, try $t=t_{\max}$ and repeatedly replace $t$ by $\zeta t$ until, for $\bm{y}=T_t(\bm{x})$ and $\bm{d}=\bm{y}-\bm{x}$,
\begin{equation}
 f_h(\bm{y})\le f_h(\bm{x})+\nabla f_h(\bm{x})^\top \bm{d}+
                       \frac{\omega}{2t}\|\bm{d}\|_2^2.
 \label{eq:comp-backtracking}
\end{equation}
A nonnegative remainder bound permits stable acceptance without subtracting nearly equal objective values. Write $d_{ij}=\bm{w}_{ij}^\top(\bm{y}-\bm{x})$ and let $m_{ij}$ be the minimum absolute residual difference along the segment from $r_{ij}(\bm{x})$ to $r_{ij}(\bm{y})$. Thus $m_{ij}=0$ when the segment crosses zero and otherwise equals the smaller endpoint magnitude. Taylor's formula gives
\begin{equation}
\begin{aligned}
&f_h(\bm y)-f_h(\bm x)-\nabla f_h(\bm x)^\top(\bm y-\bm x)
 \le\mathfrak B_h(\bm x,\bm y),\\
&\mathfrak B_h(\bm x,\bm y)=
 \frac1{2Nh}\sum_{i\ne j}\Delta_i\phi(m_{ij}/h)d_{ij}^2\\
&\hspace{38mm}\le\frac{\mathscr B_h}{2}\|\bm y-\bm x\|_2^2.
\end{aligned}
\label{eq:comp-positive-certificate}
\end{equation}
The reference implementation accepts when $\mathfrak B_h(\bm{x},\bm{y})\le\omega\|\bm{y}-\bm{x}\|_2^2/(2t)$. This sufficient test passes for $t\le\omega/\mathscr B_h$ and can be evaluated in event-row blocks.

Set $t_r=t$ and $\bm{x}^{(r+1)}=\bm{y}$. The cap $t_{\max}$ supplies a uniform step-size bound; $\zeta=\omega=1/2$ is one admissible numerical choice.

\medskip
\noindent\begin{minipage}{\textwidth}
\small
\hrule\vspace{5pt}
\textbf{Algorithm 1. Fixed-target, partially penalized proximal gradient}\par
\textit{Input:} $\bm Y,\bm\Delta,\widehat{\bm{W}},q$, fixed $h,\lambda>0$, finite $\bm{x}^{(0)}$, $t_{\max}>0$, and $\zeta,\omega\in(0,1)$.\par
\textit{Preflight:} validate dimensions, scales, event coding, PCA rank, and the nuisance-direction diagnostic in Section~5.5. Return a distinct no-event status when the event set is empty.\par
For $r=0,1,\ldots$: compute $f_h(\bm{x}^{(r)})$ and $\nabla f_h(\bm{x}^{(r)})$ by blocks; initialize $t=t_{\max}$; form~\eqref{eq:comp-prox} and backtrack until~\eqref{eq:comp-backtracking} is certified by~\eqref{eq:comp-positive-certificate}; accept the candidate; recompute its gradient and record the objective, step size, and full optimality residuals. Record the stopping status and full-coordinate residual.\par
\textit{Output:} coefficients in the original scale, frozen prediction map, all tuning values, residuals, iteration and evaluation counts, and diagnostic status. Iteration and line-search limits retain separate status codes.\par
\vspace{5pt}\hrule
\end{minipage}
\medskip

\begin{theorem}[Convergence for a fixed, attained objective]
\label{thm:comp-convergence}
Suppose the finite-sample minimizer set $\mathcal S=\arg\min Q$ is nonempty. In exact arithmetic, Algorithm~1 accepts a step after finitely many backtracking trials. If $\mathscr B_h>0$, all accepted steps satisfy
\begin{equation}
 t_r\ge t_{\min}:=\min\{t_{\max},\zeta\omega/\mathscr B_h\}>0.
 \label{eq:comp-tmin}
\end{equation}
When $\mathscr B_h=0$, set $t_{\min}=t_{\max}$. The sequence converges to a member of $\mathcal S$, and, for each $\bm{x}^\star\in\mathcal S$ and $r\ge1$,
\begin{align}
 Q(\bm{x}^{(r+1)})&\le Q(\bm{x}^{(r)})-
      \frac{2-\omega}{2t_r}\|\bm{x}^{(r+1)}-\bm{x}^{(r)}\|_2^2,
      \label{eq:comp-descent}\\
 0\le Q(\bm{x}^{(r)})-Q(\bm{x}^\star)&\le
      \frac{\|\bm{x}^{(0)}-\bm{x}^\star\|_2^2}{2t_{\min}r}.
      \label{eq:comp-rate}
\end{align}
No strong convexity or uniqueness of the minimizer is required.
\end{theorem}
For fixed $n,h,\lambda$, Theorem~\ref{thm:comp-convergence} gives an $O(r^{-1})$ objective-gap bound for Algorithm~1 in exact arithmetic. Appendix~\ref{app:computation} specializes standard proximal-gradient reasoning \citep{BeckTeboulle2009,ParikhBoyd2014} to the stated line search. Finite numerical stopping uses the residual below.

\subsection{Optimality residuals and the finite-tolerance interface}
\label{subsec:comp-stopping}
For the Karush--Kuhn--Tucker (KKT) conditions, at a finite candidate $\bm{x}=(\bm{a}^\top,\bm{\beta}^\top)^\top$, let $\bm{g}=\nabla f_h(\bm{x})$ and define
\begin{equation}
 \begin{split}
 \mathfrak r_{\rm KKT}(\bm{x})&=\max\{\|\bm{g}_a\|_\infty,\max_j\rho_j(\bm{x})\},\\
 \rho_j(\bm{x})&=
 \begin{cases}
 |g_{\beta,j}+\lambda\operatorname{sign}(\beta_j)|,&\beta_j\ne0,\\
 (|g_{\beta,j}|-\lambda)_+,&\beta_j=0.
 \end{cases}
 \end{split}
 \label{eq:comp-kkt}
\end{equation}
A maximum over an empty block is zero. This equals $\inf\{\|\bm{s}\|_\infty:\bm{s}\in\partial Q(\bm{x})\}$, and is zero if and only if the finite vector $\bm{x}$ is a global minimizer. The proximal-gradient mapping
\begin{equation}
 G_t(\bm{x})=\{\bm{x}-T_t(\bm{x})\}/t
 \label{eq:comp-mapping}
\end{equation}
provides a second, continuous fixed-point diagnostic, also zero exactly at a minimizer. At an accepted pair $(\bm{x},\bm{y})$, the explicit vector
\begin{equation}
 \begin{gathered}
\bm{s}_y=G_t(\bm{x})+\nabla f_h(\bm{y})-\nabla f_h(\bm{x})\in\partial Q(\bm{y}),\\
 \|\bm{s}_y\|_2\le(1+t\mathscr B_h)\|G_t(\bm{x})\|_2
 \end{gathered}
 \label{eq:comp-residual-vector}
\end{equation}
links the iteration to a subgradient residual. Equation~\eqref{eq:comp-kkt} uses the stored coefficient's exact zero/nonzero status.

The implementation reports absolute residuals, the declared tolerance, and the normalized residual $\mathfrak r_{\rm KKT}/\max\{\lambda,\|\bm{g}\|_\infty\}$. Numerical stationarity and nuisance-direction geometry are assessed separately. Supplementary Appendix~\ref{app:implementation} specifies the tracked arithmetic-error budget and the distinction between this numerical diagnostic and an objective-gap certificate.

\begin{corollary}[Finite optimization tolerance on the statistical event]
\label{cor:comp-statistical-tolerance}
Under the assumptions and tuning of the ordinary-model conditional theorem in Section~4, let $k=s+q$, $\bm{\vartheta}^\dagger=(\bm{\theta}_0^\top,(\widehat{\bm{B}}^\top\bm{\beta}_0)^\top,\bm{\beta}_0^\top)^\top$, and $R=8\sqrt{k}\lambda/\kappa$. On the same event used in that theorem, every finite $\widetilde{\bm{\vartheta}}$ satisfying
\begin{equation}
 \mathfrak r_{\rm KKT}(\widetilde{\bm{\vartheta}})\le\lambda/4
 \label{eq:comp-statistical-tolerance}
\end{equation}
obeys
\begin{equation}
 \begin{gathered}
\|\widetilde{\bm{\vartheta}}-\bm{\vartheta}^\dagger\|_2\le R,\\
 \|\widetilde{\bm{\beta}}-\bm{\beta}_0\|_1\le4\sqrt{k}R.
 \end{gathered}
 \label{eq:comp-statistical-transfer}
\end{equation}
Thus the conditional probability is still at least $1-2\alpha$, under the same fixed-tuning hypotheses; no extra event is introduced by this deterministic implication.
\end{corollary}
A tilted sublevel argument transfers the Section~\ref{sec:theory} score and contrast bounds to finite candidates under the same fixed-tuning conditions. The experiments use a tighter numerical threshold of $10^{-7}$; outcome-based tuning follows the separate validation procedure of Section~\ref{subsec:tuning-interface}.

\subsection{Unpenalized escape directions and finite existence}
\label{subsec:comp-existence}
Define the event-anchored nuisance differences $\bm{a}_{ij}=\widehat{\bm{A}}_i-\widehat{\bm{A}}_j$ and the flat-direction space
\[
 \mathcal N_A=\{\bm{v}\in\mathbb R^q:\bm{a}_{ij}^\top \bm{v}=0
                         \text{ whenever }\Delta_i=1,\ j\ne i\}.
\]
A \emph{strict escape direction} is a vector $\bm{v}$ such that
\begin{equation}
 \begin{gathered}
\bm{a}_{ij}^\top \bm{v}\le0\quad(\Delta_i=1,\ j\ne i),
 \\
 \bm{a}_{ij}^\top \bm{v}<0\quad\text{for at least one such pair}.
 \end{gathered}
 \label{eq:comp-escape}
\end{equation}
Strictness for just one pair suffices because $\ell_h'(r)>0$ at every finite $r$ when $h>0$.

\begin{proposition}[An exact-arithmetic finite-existence criterion]
\label{prop:comp-existence}
For finite $\bm Y,\widehat{\bm{W}}$, $h,\lambda>0$, the set $\arg\min Q$ is nonempty if and only if no strict escape direction in~\eqref{eq:comp-escape} exists. If no such direction exists, every nonempty sublevel set is compact after projecting $\bm{a}$ onto $\mathcal N_A^\perp$. Directions in $\mathcal N_A$ leave the objective unchanged and need not imply nonexistence. When $q=0$, a finite minimum always exists.
\end{proposition}
Proposition~\ref{prop:comp-existence} characterizes finite attainment of the fixed criterion, including flat nuisance directions. The coefficient guarantees in Section~\ref{sec:theory} additionally impose restricted-design separation.

If $m>0$, choose an event index $i_0$ and let $\mathcal E=\{i:\Delta_i=1\}$, $\mathcal C=\{j:\Delta_j=0\}$. The escape test can be expressed with $O(nq)$ coefficients rather than all pair differences: find $\bm{v}\in\mathbb R^q$ satisfying
\begin{equation}
 \begin{aligned}
 (\widehat{\bm{A}}_i-\widehat{\bm{A}}_{i_0})^\top \bm{v}&=0&& (i\in\mathcal E),\\
 (\widehat{\bm{A}}_j-\widehat{\bm{A}}_{i_0})^\top \bm{v}&\ge0&& (j\in\mathcal C),\\
 \sum_{j\in\mathcal C}(\widehat{\bm{A}}_j-\widehat{\bm{A}}_{i_0})^\top \bm{v}&\ge1.&
 \end{aligned}
 \label{eq:comp-escape-lp}
\end{equation}
The last constraint fixes the scale of a positive strict margin. In exact arithmetic, feasibility is equivalent to~\eqref{eq:comp-escape}. Numerical output includes equality residuals, inequality margins, and an explicit inconclusive status when the margins cannot be resolved.

For $m=0$, $Q=\lambda\|\bm{\beta}\|_1$ and every $\bm{a}$ with $\bm{\beta}=0$ is optimal, but the sample provides no event information. For $m=n$, strict escape is impossible because both directions of every comparison are present. Rank deficiency can still leave flat nuisance coordinates. These cases distinguish nonexistence, nonuniqueness, and lack of event information.

\subsection{Regularization paths and stable arithmetic}
\label{subsec:comp-paths}
For fixed $(K,h)$, suppose $\bm{a}^{(0)}$ attains the minimum of $\bm{a}\mapsto f_h(\bm{a},0)$. The level
\begin{equation}
 \lambda_{\rm zero}=\|\nabla_\beta f_h(\bm{a}^{(0)},0)\|_\infty
 \label{eq:comp-lambda-zero}
\end{equation}
makes $(\bm{a}^{(0)},0)$ a global minimizer for every $\lambda\ge\lambda_{\rm zero}$ by the full KKT conditions. The threshold is evaluated after checking nuisance-only stationarity. Penalty paths retain strictly positive levels and use the preceding successful fit as a warm start. This construction follows regularized AFT practice \citep{SuderMolstad2022}, with the null fit including the unpenalized block.

For a very negative residual difference, evaluating $r\Phi(r/h)+h\phi(r/h)$ naively subtracts nearly equal quantities. A useful stable identity, with $z=|r|/h$, is
\begin{equation}
 \ell_h(r)=r_++h\phi(z)J(z),\qquad
 J(z)=\int_0^\infty u\exp(-zu-u^2/2)\,\mathrm{d}u>0.
 \label{eq:comp-stable-loss}
\end{equation}
For $z>0$, an alternating tail expansion is
\begin{equation}
 \begin{split}
 J(z)&=\sum_{j=0}^{M-1}\frac{(-1)^j(2j+1)!!}{z^{2j+2}}+\mathcal E_M(z),\\
 |\mathcal E_M(z)|&\le\frac{(2M+1)!!}{z^{2M+2}}.
 \end{split}
 \label{eq:comp-tail-expansion}
\end{equation}
Here $M$ is the expansion length. The implementation combines normal-tail routines, log-scale evaluation, and the explicit remainder bound. Numerical metadata track tail approximation, unresolved arithmetic, and the line-search budget; details appear in Supplementary Appendix~\ref{app:implementation}.

Every candidate retains its declared objective, complete-coordinate stopping rule, and frozen prediction map. Output states distinguish stationarity to tolerance, iteration limits, numerical or line-search failures, nuisance escape, flat directions, and event-free samples.
